% ============================================================
\chapter{Introduction aux Cha\^ines de Markov}
\label{ch:markov}
% ============================================================

\section{Introduction}

Les cha\^ines de Markov\index{cha\^ine de Markov} constituent l'un des mod\`eles
stochastiques les plus fondamentaux. Elles mod\'elisent l'\'evolution al\'eatoire d'un
syst\`eme dont le futur ne d\'epend du pass\'e qu'\`a travers le pr\'esent (propri\'et\'e
de Markov). Leurs applications couvrent la physique (marches al\'eatoires), la biologie
(\'evolution de populations), l'informatique (algorithmes de Monte-Carlo, PageRank) et
la finance.

\begin{intuition}[title=\textbf{La propri\'et\'e de Markov~: absence de m\'emoire}]
Imaginez un pion sur un \'echiquier qui se d\'eplace al\'eatoirement. \`A chaque \'etape,
la probabilit\'e de se d\'eplacer vers une case voisine ne d\'epend que de la position
actuelle du pion, pas de la trajectoire pass\'ee. C'est la propri\'et\'e de Markov~:
\og le futur est ind\'ependant du pass\'e, conditionn\'e au pr\'esent \fg.
\end{intuition}

% ------------------------------------------------------------
\section{D\'efinitions}
% ------------------------------------------------------------

\begin{definition}[Cha\^ine de Markov]
\label{def:chaine_markov}
\index{cha\^ine de Markov}
Soit $E$ un ensemble d\'enombrable (l'\emph{espace d'\'etats}). Un processus
stochastique $(X_n)_{n \geq 0}$ \`a valeurs dans $E$ est une \emph{cha\^ine de Markov}
si pour tous $n \geq 0$ et tous \'etats $i_0, \ldots, i_{n-1}, i, j \in E$~:
\[
    \PP(X_{n+1} = j \mid X_n = i, X_{n-1} = i_{n-1}, \ldots, X_0 = i_0)
    = \PP(X_{n+1} = j \mid X_n = i),
\]
d\`es que le conditionnement est bien d\'efini ($\PP(X_n = i, \ldots) > 0$).
\end{definition}

\begin{definition}[Cha\^ine homog\`ene et matrice de transition]
\label{def:matrice_transition}
\index{matrice de transition}
La cha\^ine est dite \emph{homog\`ene} (ou \emph{stationnaire}) si les probabilit\'es
de transition ne d\'ependent pas de $n$~:
\[
    p_{ij} = \PP(X_{n+1} = j \mid X_n = i), \qquad \forall n \geq 0.
\]
La matrice $P = (p_{ij})_{i,j \in E}$ est appel\'ee \emph{matrice de transition}.
Elle v\'erifie~:
\begin{enumerate}
    \item[(i)] $p_{ij} \geq 0$ pour tous $i, j$\,;
    \item[(ii)] $\sum_{j \in E} p_{ij} = 1$ pour tout $i$ \quad (matrice stochastique).
\end{enumerate}
\end{definition}

\begin{definition}[Loi initiale]
La \emph{loi initiale} est la loi de $X_0$, c'est-\`a-dire le vecteur
$\mu_0 = (\PP(X_0 = i))_{i \in E}$.
La loi de la cha\^ine est enti\`erement d\'etermin\'ee par $\mu_0$ et $P$.
\end{definition}

\begin{proposition}[Loi jointe]
Pour tous $i_0, i_1, \ldots, i_n \in E$~:
\[
    \PP(X_0 = i_0, X_1 = i_1, \ldots, X_n = i_n)
    = \mu_0(i_0)\, p_{i_0 i_1}\, p_{i_1 i_2} \cdots p_{i_{n-1} i_n}.
\]
\end{proposition}

\begin{exemple}[Marche al\'eatoire sur $\Z$]
\index{marche al\'eatoire}
Soit $X_0 = 0$ et $X_{n+1} = X_n + \xi_{n+1}$ o\`u les $(\xi_n)$ sont i.i.d.\ avec
$\PP(\xi = +1) = p$ et $\PP(\xi = -1) = 1-p$. C'est une cha\^ine de Markov sur
$E = \Z$ avec $p_{i,i+1} = p$ et $p_{i,i-1} = 1-p$.
\end{exemple}

\begin{exemple}[Mod\`ele d'Ehrenfest]
$2N$ particules sont r\'eparties entre deux bo\^ites. \`A chaque \'etape, on choisit une
particule au hasard et on la d\'eplace dans l'autre bo\^ite. Si $X_n$ est le nombre de
particules dans la premi\`ere bo\^ite, $(X_n)$ est une cha\^ine de Markov sur
$\{0, 1, \ldots, 2N\}$ avec $p_{k,k-1} = k/(2N)$ et $p_{k,k+1} = 1 - k/(2N)$.
\end{exemple}

% ------------------------------------------------------------
\section{Probabilit\'es de transition en $n$ pas}
% ------------------------------------------------------------

\begin{definition}
Les \emph{probabilit\'es de transition en $n$ pas} sont~:
\[
    p_{ij}^{(n)} = \PP(X_n = j \mid X_0 = i).
\]
\end{definition}

\begin{theoreme}[\'Equation de Chapman-Kolmogorov]
\label{thm:chapman_kolmogorov}
\index{\'equation de Chapman-Kolmogorov}
Pour tous $m, n \geq 0$ et tous $i, j \in E$~:
\[
    p_{ij}^{(m+n)} = \sum_{k \in E} p_{ik}^{(m)}\, p_{kj}^{(n)}.
\]
En notation matricielle~: $P^{(m+n)} = P^m \cdot P^n$, soit $P^{(n)} = P^n$.
\end{theoreme}

\begin{proof}
Par la formule des probabilit\'es totales~:
\begin{align*}
    p_{ij}^{(m+n)} &= \PP(X_{m+n} = j \mid X_0 = i)
    = \sum_{k \in E} \PP(X_{m+n} = j \mid X_m = k) \PP(X_m = k \mid X_0 = i) \\
    &= \sum_{k \in E} p_{kj}^{(n)}\, p_{ik}^{(m)}.
\end{align*}
\end{proof}

% ------------------------------------------------------------
\section{Classification des \'etats}
% ------------------------------------------------------------

\begin{definition}[Accessibilit\'e et communication]
\index{communication}
\begin{itemize}
    \item L'\'etat $j$ est \emph{accessible} depuis $i$, not\'e $i \to j$, s'il
    existe $n \geq 0$ tel que $p_{ij}^{(n)} > 0$.
    \item Deux \'etats $i$ et $j$ \emph{communiquent}, not\'e $i \leftrightarrow j$,
    si $i \to j$ et $j \to i$.
\end{itemize}
La relation $\leftrightarrow$ est une relation d'\'equivalence. Ses classes sont les
\emph{classes de communication}.
\end{definition}

\begin{definition}[Cha\^ine irr\'eductible]
\index{irr\'eductible}
La cha\^ine est \emph{irr\'eductible} si tous les \'etats communiquent, c'est-\`a-dire
s'il n'y a qu'une seule classe de communication.
\end{definition}

\begin{definition}[Temps de retour et r\'ecurrence]
\index{r\'ecurrence}\index{transience}
Le \emph{temps de premier retour} en $i$ est~:
$T_i = \inf\{n \geq 1 : X_n = i\}$.
\begin{itemize}
    \item L'\'etat $i$ est \emph{r\'ecurrent} si $\PP_i(T_i < \infty) = 1$.
    \item L'\'etat $i$ est \emph{transitoire} si $\PP_i(T_i < \infty) < 1$.
\end{itemize}
\end{definition}

\begin{theoreme}[Caract\'erisation de la r\'ecurrence]
\label{thm:recurrence}
Un \'etat $i$ est r\'ecurrent si et seulement si $\sum_{n=0}^{\infty} p_{ii}^{(n)} = +\infty$.
Un \'etat $i$ est transitoire si et seulement si $\sum_{n=0}^{\infty} p_{ii}^{(n)} < +\infty$.
\end{theoreme}

\begin{proof}
Soit $N_i = \sum_{n=1}^{\infty} \ind_{\{X_n = i\}}$ le nombre de retours en $i$.
On a $\E_i[N_i] = \sum_{n=1}^{\infty} p_{ii}^{(n)}$. De plus,
$\PP_i(N_i \geq k) = \PP_i(T_i < \infty)^k$, donc $N_i$ suit une loi g\'eom\'etrique.
Si $i$ est r\'ecurrent, $\PP_i(T_i < \infty) = 1$, donc $\E_i[N_i] = +\infty$.
Si $i$ est transitoire, $\PP_i(T_i < \infty) < 1$, donc
$\E_i[N_i] = \PP_i(T_i < \infty)/(1 - \PP_i(T_i < \infty)) < +\infty$.
\end{proof}

\begin{proposition}[Propri\'et\'e de classe]
La r\'ecurrence et la transience sont des propri\'et\'es de classe~: si $i$
communique avec $j$, alors $i$ est r\'ecurrent si et seulement si $j$ l'est.
\end{proposition}

\begin{definition}[\'Etat absorbant]
\index{\'etat absorbant}
Un \'etat $i$ est \emph{absorbant} si $p_{ii} = 1$. Un \'etat absorbant est
r\'ecurrent.
\end{definition}

\begin{definition}[P\'eriode]
\index{p\'eriode}
La \emph{p\'eriode} d'un \'etat $i$ est $d(i) = \gcd\{n \geq 1 : p_{ii}^{(n)} > 0\}$.
L'\'etat $i$ est \emph{ap\'eriodique} si $d(i) = 1$. La p\'eriode est une propri\'et\'e
de classe.
\end{definition}

% ------------------------------------------------------------
\section{Distributions stationnaires}
% ------------------------------------------------------------

\begin{definition}[Distribution stationnaire]
\label{def:stationnaire}
\index{distribution stationnaire}
Un vecteur de probabilit\'e $\pi = (\pi_i)_{i \in E}$ (avec $\pi_i \geq 0$ et
$\sum_i \pi_i = 1$) est une \emph{distribution stationnaire} (ou \emph{invariante})
pour la matrice $P$ si~:
\[
    \pi P = \pi, \qquad \text{c'est-\`a-dire} \quad
    \pi_j = \sum_{i \in E} \pi_i\, p_{ij}, \quad \forall j \in E.
\]
\end{definition}

\begin{theoreme}[Existence pour les cha\^ines irr\'eductibles r\'ecurrentes positives]
\label{thm:existence_stationnaire}
Soit $(X_n)$ une cha\^ine de Markov irr\'eductible et r\'ecurrente. Alors~:
\begin{enumerate}
    \item Il existe une mesure invariante $\pi$ (unique \`a un facteur multiplicatif
    pr\`es) donn\'ee par $\pi_j = 1/\E_j[T_j]$.
    \item $\pi$ est une distribution de probabilit\'e (i.e.\ $\sum_j \pi_j = 1$) si
    et seulement si la cha\^ine est \emph{r\'ecurrente positive}, c'est-\`a-dire
    $\E_i[T_i] < \infty$ pour tout (ou de fa\c{c}on \'equivalente, pour un) \'etat $i$.
\end{enumerate}
\end{theoreme}

\begin{remarque}
Si l'espace d'\'etats $E$ est fini et la cha\^ine est irr\'eductible, alors elle est
automatiquement r\'ecurrente positive et admet une unique distribution stationnaire.
\end{remarque}

% ------------------------------------------------------------
\section{Th\'eor\`eme ergodique}
% ------------------------------------------------------------

\begin{theoreme}[Th\'eor\`eme ergodique pour les cha\^ines de Markov]
\label{thm:ergodique}
\index{th\'eor\`eme ergodique}
Soit $(X_n)$ une cha\^ine de Markov irr\'eductible, r\'ecurrente positive et
ap\'eriodique, de distribution stationnaire $\pi$. Alors pour tout \'etat initial
$i$ et tout \'etat $j$~:
\[
    \lim_{n \to \infty} p_{ij}^{(n)} = \pi_j.
\]
De plus, pour toute fonction born\'ee $f : E \to \R$~:
\[
    \frac{1}{n} \sum_{k=0}^{n-1} f(X_k) \xrightarrow[n \to \infty]{\text{p.s.}} \sum_{j \in E} f(j)\, \pi_j.
\]
\end{theoreme}

\begin{formules}[title=\textbf{Formules cl\'es des cha\^ines de Markov}]
\begin{itemize}
    \item \textbf{Chapman-Kolmogorov~:} $P^{(m+n)} = P^m \cdot P^n$
    \item \textbf{Distribution au temps $n$~:} $\mu_n = \mu_0 P^n$
    \item \textbf{Stationnarit\'e~:} $\pi P = \pi$
    \item \textbf{Lien avec le temps de retour~:} $\pi_j = 1/\E_j[T_j]$
    \item \textbf{Convergence~:} $\lim_{n \to \infty} P^n = \mathbf{1} \pi^T$
    (si irr\'eductible, r\'ecurrente positive, ap\'eriodique)
\end{itemize}
\end{formules}

% ------------------------------------------------------------
\section{R\'eversibilit\'e et \'equations de bilan d\'etaill\'e}
% ------------------------------------------------------------

\begin{definition}[R\'eversibilit\'e]
\index{r\'eversibilit\'e}\index{bilan d\'etaill\'e}
Une cha\^ine de Markov de matrice de transition $P$ et de distribution stationnaire
$\pi$ est dite \emph{r\'eversible} si les \emph{\'equations de bilan d\'etaill\'e}
sont satisfaites~:
\[
    \pi_i\, p_{ij} = \pi_j\, p_{ji}, \qquad \forall i, j \in E.
\]
\end{definition}

\begin{proposition}
Si $\pi$ v\'erifie les \'equations de bilan d\'etaill\'e, alors $\pi$ est une
distribution stationnaire.
\end{proposition}

\begin{proof}
En sommant sur $i$~:
$\sum_i \pi_i\, p_{ij} = \sum_i \pi_j\, p_{ji} = \pi_j \sum_i p_{ji} = \pi_j$,
car $\sum_i p_{ji} = 1$.
\end{proof}

\begin{exemple}[Marche al\'eatoire sur un graphe]
Soit $G = (V, E)$ un graphe connexe fini non orient\'e. La marche al\'eatoire simple
sur $G$ a pour matrice de transition $p_{ij} = 1/\mathrm{deg}(i)$ si $(i,j)$ est une
ar\^ete. La distribution stationnaire est $\pi_i = \mathrm{deg}(i)/(2\abs{E})$, et
les \'equations de bilan d\'etaill\'e sont satisfaites.
\end{exemple}

\begin{attention}[title=\textbf{Bilan d\'etaill\'e $\neq$ stationnarit\'e}]
Le bilan d\'etaill\'e est une condition \emph{suffisante} mais pas n\'ecessaire pour
la stationnarit\'e. Il existe des cha\^ines de Markov avec une distribution
stationnaire qui ne v\'erifie pas le bilan d\'etaill\'e (par exemple, les cha\^ines
avec un cycle orient\'e).
\end{attention}

% ------------------------------------------------------------
\section{Exercices}
% ------------------------------------------------------------

\begin{exercice}
Soit $(X_n)$ la marche al\'eatoire sym\'etrique sur $\Z$ ($p = 1/2$).
Montrer que chaque \'etat est r\'ecurrent.
\emph{Indication}~: utiliser $p_{00}^{(2n)} = \binom{2n}{n} 2^{-2n} \sim (\pi n)^{-1/2}$.
\end{exercice}

\begin{exercice}
Montrer que la marche al\'eatoire sym\'etrique sur $\Z^3$ est transitoire.
\emph{Indication}~: montrer que $\sum_n p_{00}^{(n)} < \infty$.
\end{exercice}

\begin{exercice}
Soit $E = \{1, 2, 3\}$ et
$P = \begin{pmatrix} 0 & 1/2 & 1/2 \\ 1/3 & 1/3 & 1/3 \\ 1/2 & 0 & 1/2 \end{pmatrix}$.
\begin{enumerate}
    \item V\'erifier que la cha\^ine est irr\'eductible et ap\'eriodique.
    \item Calculer la distribution stationnaire $\pi$.
    \item La cha\^ine est-elle r\'eversible~?
\end{enumerate}
\end{exercice}

\begin{exercice}[Mod\`ele de file d'attente]
\index{file d'attente}
Soit $N \geq 1$ la capacit\'e d'une file d'attente. \`A chaque \'etape, un client arrive
avec probabilit\'e $\lambda$ et un client est servi avec probabilit\'e $\mu$ (si la file
n'est pas vide), ind\'ependamment. Soit $X_n$ le nombre de clients au temps $n$.
Montrer que $(X_n)$ est une cha\^ine de Markov, \'ecrire sa matrice de transition,
et trouver sa distribution stationnaire lorsque $\lambda < \mu$.
\end{exercice}

\begin{exercice}[Algorithme de Metropolis-Hastings]
\index{Metropolis-Hastings}
Soit $\pi$ une distribution de probabilit\'e cible sur un ensemble fini $E$ et soit $Q$
une matrice de transition irr\'eductible (la \emph{proposition}). On d\'efinit la matrice
de transition~:
\[
    p_{ij} = q_{ij} \min\!\left(1, \frac{\pi_j q_{ji}}{\pi_i q_{ij}}\right)
    \quad (i \neq j), \qquad
    p_{ii} = 1 - \sum_{j \neq i} p_{ij}.
\]
Montrer que $\pi$ est une distribution stationnaire pour $P$ et que le bilan
d\'etaill\'e est satisfait.
\end{exercice}

\begin{exercice}
Soit $(X_n)$ une cha\^ine de Markov irr\'eductible sur un espace d'\'etats fini
$E$ avec $\abs{E} = N$. Montrer qu'il existe un unique vecteur stationnaire $\pi$
et que $\pi_i > 0$ pour tout $i$.
\end{exercice}
